
\chapter{Supplementary Prompts and Case Studies for
Simulation-Grounded Architectural Explanation}
\label{app:p5_supplementary}

This appendix provides supplementary
implementation details and recorded
case studies for the simulation-grounded
architectural explanation framework
presented in Part B of
Chapter~\ref{chap:your_chapter_label}.

Unlike Part A, which examines
LLM reasoning during DSE,
Part B investigates the correctness
of architectural explanations
generated in response to user questions.

Four evidence-access conditions
are examined: Dynamic Memory (C1),
Pure Graph (C2), Evidence-Enriched
Graph (C3), and Question-Focused
GraphRAG (C4).

The examples illustrate how
simulation evidence accumulated
during DSE is represented,
organized, enriched, and retrieved
to support architectural explanations.

The appendix also presents selected
LLM prompt instructions and recorded
verification case studies to clarify
how simulation evidence is supplied
to the LLM and how the resulting
explanations are evaluated.

The case studies are drawn from
recorded experimental outputs and
are intended to illustrate the
operation and limitations of
individual mechanisms.

Their aggregate effectiveness
is evaluated separately in the
main Experimental Results section.



\section{LLM Prompt Templates}
\label{app:p5_prompts}

This section presents selected
prompt instructions used to generate
simulation-grounded architectural
explanations in Part B.

Unlike the Reflection and ExpeL
prompts documented in
Appendix~\ref{app:p4_supplementary},
which support reasoning improvement
during DSE, the Part B prompts
guide the LLM in answering user
questions using previously
accumulated architectural evidence.

The prompt instructions emphasize
the use of recorded simulation
results, accurate interpretation
of architectural observations,
and appropriate communication
of evidence limitations.

Only the instructions relevant
to architectural explanation
are presented here to avoid
repeating shared prompt components
across the four evidence conditions.


\subsection{Architectural Explanation Prompt}
\label{app:p5_explanation_prompt}

The Runtime Knowledge-Grounded Tutor
uses a shared system prompt to generate
architectural explanations from evidence
accumulated during the current DSE run.

The prompt establishes a clear boundary
between general architectural knowledge
and measured runtime evidence.

It also requires the Tutor to distinguish
direct observations from interpretations,
avoid unsupported causal conclusions,
and explicitly acknowledge insufficient
simulation evidence.

These instructions are shared across
the four evidence-access conditions
(C1--C4), ensuring that the experimental
comparison focuses on differences in
the supplied evidence rather than
differences in the Tutor's general
explanation policy.

Listing~\ref{lst:app_p5_explanation_prompt}
presents selected instructions from
the implemented Tutor system prompt.

\begin{lstlisting}[
    style=remodeljson,
    caption={Selected instructions from the
    Runtime Knowledge-Grounded Tutor prompt.},
    label={lst:app_p5_explanation_prompt},
    breaklines=true,
    basicstyle=\ttfamily\footnotesize
]
RUNTIME KNOWLEDGE-GROUNDED TUTOR

Your purpose is to help students understand
the current optimization run using only
the runtime knowledge supplied in
STUDENT_CONTEXT.

KNOWLEDGE BOUNDARY

General educational knowledge may be used
only to explain computer architecture or
Design Space Exploration concepts.

General educational knowledge must never
be presented as runtime evidence.

Runtime conclusions, observations,
comparisons, recommendations, and
measured results must be supported
only by STUDENT_CONTEXT.

When the runtime evidence is insufficient,
state:

"The current optimization run does not
contain enough evidence to answer this
with confidence."

EVIDENCE TYPES

Always distinguish among:

- OBSERVED FACT
- DIRECT COMPARISON
- RUNTIME RELATIONSHIP
- VERIFIED REASONING
- REFLECTION
- LESSON
- UNVERIFIED HYPOTHESIS

CAUSALITY RULES

When multiple parameters changed,
never begin by stating that the metric
changed "because of" one parameter.

First identify all parameters that
changed between the compared iterations.

Then determine whether the evidence
contains an exact parameter-isolated
comparison.

If the parameter was not isolated,
describe only an association,
consistency with a claim, or
possible contribution.

Do not claim that one parameter
caused a metric change unless
the supplied evidence indicates
sufficient parameter isolation.

STUDENT QUESTIONS

Answer every part that is supported
by the current-run evidence.

Do not force the question into
one category.
\end{lstlisting}

The excerpt highlights the main
evidence-grounding requirements
used during architectural question
answering.

In particular, the Tutor must
support run-specific statements
using supplied simulation evidence
and avoid presenting previously
generated claims, reflections,
or lessons as independently
measured facts.

The causality instructions also
prevent the Tutor from attributing
observed metric changes to a
single architectural parameter
when other parameters changed
simultaneously.

These requirements provide a
common explanation policy for
C1--C4, while the evidence
representation and retrieval
mechanisms determine which
runtime information is supplied
to the Tutor.


\paragraph{Response Format and Evidence Traceability.}

In addition to the evidence-grounding
instructions, the Tutor follows a
structured response format that
separates the user's question,
the generated explanation, and
its supporting runtime evidence.

Listing~\ref{lst:app_p5_response_format}
presents the relevant instructions
from the implemented Tutor prompt.

\begin{lstlisting}[
    style=remodeljson,
    caption={Tutor response format and
    evidence-traceability instructions.},
    label={lst:app_p5_response_format},
    breaklines=true,
    basicstyle=\ttfamily\footnotesize
]
RESPONSE FORMAT

Always return the answer using
exactly this structure:

Q:
Repeat the student's complete question.

A:
Answer the question clearly using only
the supplied current-run evidence.

Evidence:
Show the specific runtime information
that supports the answer.

SOURCE TRACEABILITY

The Evidence section must list only
the runtime evidence actually used
to answer the student's question.

For each source, identify:
- the iteration number;
- the evidence type.

Do not cite an iteration unless
information from that iteration
was actually used.

Do not invent file names,
iteration numbers, confidence values,
evidence counts, or verification outcomes.
\end{lstlisting}

The required response format makes
the Tutor's explanation and its
supporting evidence explicit.

The \texttt{Q} field identifies
the architectural question,
the \texttt{A} field contains
the generated explanation,
and the \texttt{Evidence} field
provides the relevant runtime
observations supporting that answer.

The traceability instructions
require the Tutor to identify
the source iterations and
evidence types used in its
explanation.

These requirements are intended
to improve transparency and
reduce unsupported architectural
statements.

However, following a structured
response format does not guarantee
factual correctness.

Part B therefore independently
evaluates the generated explanations
against authoritative simulation
evidence using human and
LLM-based verification.



\section{From Dynamic Memory to Pure Graph}
\label{app:p5_c1_c2_example}

This case study illustrates how the same
simulation-grounded architectural evidence
is represented under C1 Dynamic Memory
and C2 Pure Graph.

The example is drawn from a recorded
Matrix Multiplication DSE execution
and examines an architectural prediction
concerning L1D cache size and area.

The objective is to demonstrate what
information is available to the Tutor
under each condition, how C2 improves
the organization of that information,
and what limitations remain.

Importantly, C2 does not introduce
new simulation observations or
accumulated relationship assessments.
Instead, it reorganizes the evidence
already preserved in C1.

\subsection{C1: Chronological Runtime Evidence}

Under C1, executed configurations,
simulation measurements, architectural
claims, and their supporting or
contradicting evidence are stored
within chronological Dynamic Memory
experiences.

During Iteration 2, the LLM
generated the prediction:

\begin{quote}
``Reducing \texttt{l1d\_size}
may lower area.''
\end{quote}

Listing~\ref{lst:app_b_c1_memory}
presents selected fields from the
recorded Iteration 2 experience.

\begin{lstlisting}[
    style=remodeljson,
    caption={Selected architectural information
    preserved in C1 Dynamic Memory.},
    label={lst:app_b_c1_memory}
]
{
  "iteration": 2,
  "design": {
    "l1d_size": "8kB",
    "rob_entries": 64
  },
  "simulation_results": {
    "area": 4.2549
  },
  "extracted_claims": [
    {
      "claim_text":
        "Reducing l1d_size may lower area.",
      "parameter_direction": "decrease",
      "expected_direction": "decrease"
    }
  ]
}
\end{lstlisting}

The corresponding runtime evidence
compares the executed configurations
from Iterations 1 and 2.

Table~\ref{tab:app_b_c1_observation}
summarizes the relevant measurements.

\begin{table}[htbp]
\centering
\caption{Recorded architectural observations
from Iterations 1 and 2.}
\label{tab:app_b_c1_observation}

\renewcommand{\arraystretch}{1.2}
\small

\begin{tabular}{lcc}
\toprule
\textbf{Parameter / Metric} &
\textbf{Iteration 1} &
\textbf{Iteration 2} \\
\midrule

L1D cache size & 16kB & 8kB \\
ROB entries & 128 & 64 \\
Area & 4.3912 & 4.2549 \\

\bottomrule
\end{tabular}
\end{table}

The observations show that
L1D cache size decreased
from 16kB to 8kB, while
area decreased from 4.3912
to 4.2549.

This agrees with the predicted
direction of the relationship.

However, ROB entries also decreased
from 128 to 64. Therefore,
the observed area reduction
cannot be attributed solely
to L1D cache size.

C1 preserves the information
needed to identify this
confounded comparison.

\paragraph{Strength and Limitation of C1.}

The main strength of C1 is that
it preserves the original
simulation-grounded experiences,
including executed designs,
measurements, claims, and
runtime evidence.

However, information about
the same architectural relationship
may appear in several different
iteration experiences.

For example, to determine whether
reducing L1D cache size consistently
reduces area across the completed
run, the Tutor must locate
relevant observations from
multiple iterations and
combine their evidence.

Thus, C1 provides the underlying
information but does not organize
related observations into
explicit graph connections
or accumulated relationship
assessments.

\subsection{C2: Organizing Evidence as a Pure Graph}

C2 addresses the organizational
limitation of C1 by transforming
the same Dynamic Memory evidence
into interconnected graph nodes
and edges.

Instead of relying only on
chronological iteration records,
C2 explicitly connects
iterations, executed designs,
architectural parameters,
metrics, claims, and evidence.

For the same Iteration 2 example,
Listing~\ref{lst:app_b_c2_edges}
presents selected graph connections.

\begin{lstlisting}[
    style=remodeljson,
    caption={Selected C2 graph connections
    for the same architectural observation.},
    label={lst:app_b_c2_edges}
]
[
  {
    "type": "has_design",
    "source": "i2",
    "target": "d2"
  },
  {
    "type": "has_parameter",
    "source": "d2",
    "target": "p::l1d_size",
    "properties": {
      "value": "8kB"
    }
  },
  {
    "type": "has_metric",
    "source": "d2",
    "target": "m::area",
    "properties": {
      "value": 4.2549
    }
  },
  {
    "type": "has_claim",
    "source": "i2",
    "target": "c::2::0"
  }
]
\end{lstlisting}

The \texttt{has\_design} edge
connects Iteration 2 to its
executed architectural configuration.

The \texttt{has\_parameter}
and \texttt{has\_metric}
edges associate that configuration
with its L1D cache size
and measured area.

The \texttt{has\_claim} edge
connects the iteration to
its recorded architectural
prediction.

Additional graph connections
associate the claim with
its corresponding parameter,
metric, and runtime evidence.

Therefore, C2 makes the
relationships among recorded
architectural entities explicit
without changing the underlying
simulation observations.

\paragraph{What C2 Improves over C1.}

Under C1, the Tutor encounters
architectural information within
chronological iteration experiences.

Under C2, the same information
is organized through explicit
connections among the relevant
architectural entities.

For example, the Tutor can
identify that Iteration 2
contains an executed design
with \texttt{l1d\_size = 8kB},
a measured area of 4.2549,
and an associated architectural
claim about L1D cache size.

This organization makes the
connections between the claim,
the executed configuration,
and the measured outcome
explicit in the representation.

The improvement is therefore
structural rather than empirical:
C2 organizes existing evidence
but does not introduce
additional simulation knowledge.

\paragraph{Remaining Limitation of C2.}

Although C2 connects individual
claims to their runtime evidence,
it does not explicitly provide
an accumulated assessment of
the same architectural relationship
across the completed DSE run.

For example, the graph preserves
the Iteration 2 observation
supporting the predicted decrease
in area.

However, answering whether this
relationship is supported
across multiple iterations
requires examining other
relevant claim-evidence records
and combining their outcomes.

The Tutor must determine
which iterations support
or contradict the relationship,
whether observations are
confounded, and what causal
conclusion is justified.

Thus, C2 improves evidence
organization but does not
eliminate the need to aggregate
individual observations.

\paragraph{Summary and Motivation for C3.}

This example demonstrates
the distinction between
preserving evidence and
organizing evidence.

C1 provides chronological
simulation-grounded experiences,
while C2 represents the same
information through connected
graph entities.

The Pure Graph transformation
was validated across 27 held-out
DSE runs, all of which passed
the implemented semantic
recoverability checks.

However, neither C1 nor C2
explicitly provides the
accumulated relationship
assessment required to answer
questions about evidence
collected across multiple
iterations.

This remaining limitation
motivates C3 Evidence-Enriched
Graph, which adds deterministic
relationship-level assessments
derived from the recorded
runtime evidence.




% =====================================================
% APPENDIX — FROM PURE GRAPH TO ENRICHED GRAPH
% =====================================================

\section{From Pure Graph to Evidence-Enriched Graph}
\label{app:p5_c2_c3_comparison}
\label{app:p5_c3_example}

This section continues the same
Matrix Multiplication example
introduced in
Section~\ref{app:p5_c1_c2_example}.

The architectural prediction is:

\begin{quote}
``Reducing \texttt{l1d\_size}
may lower area.''
\end{quote}

The previous section showed that
C1 preserves the runtime observations,
while C2 organizes those observations
into connected graph nodes and edges.

However, C2 does not explicitly
summarize what all the recorded
observations indicate about
the architectural relationship.

C3 Evidence-Enriched Graph
addresses this limitation by adding
accumulated relationship assessments
derived from the same runtime evidence.

The purpose of this example is
to show what additional information
C3 provides, how it helps answer
architectural questions, and
what limitation remains.

% =====================================================
% C2 — INDIVIDUAL EVIDENCE
% =====================================================

\subsection{C2: Individual Evidence Records}

As shown in the previous section,
C2 connects architectural claims
to their corresponding runtime evidence.

For example, during Iteration 2,
L1D cache size decreased from
16kB to 8kB, while area decreased
from 4.3912 to 4.2549.

The observed change supports
the predicted direction.

However, ROB entries also decreased
from 128 to 64, so the comparison
does not isolate the effect
of L1D cache size.

C2 preserves this information
within the evidence node
associated with the claim.

Listing~\ref{lst:app_p5_c2_evidence}
shows selected fields from
that evidence record.

\begin{lstlisting}[
    style=remodeljson,
    caption={Individual runtime evidence
    available in C2 Pure Graph.},
    label={lst:app_p5_c2_evidence}
]
{
  "supporting_evidence": [
    {
      "comparison_iteration": 1,
      "current_iteration": 2,
      "parameter": "l1d_size",
      "previous_parameter_value": "16kB",
      "current_parameter_value": "8kB",
      "metric": "area",
      "previous_metric_value": 4.3912,
      "current_metric_value": 4.2549,
      "evidence_category": "supporting",
      "exact_parameter_isolation": false
    }
  ]
}
\end{lstlisting}

This record tells the Tutor
that the observation supports
the predicted direction,
but does not establish
a single-parameter cause.

\paragraph{What Is Still Missing in C2?}

The Iteration 2 observation
is only one part of the
completed DSE run.

Other iterations also contain
claims and evidence about
the relationship between
L1D cache size and area.

Suppose the student asks:

\begin{quote}
``Based on all the evidence
collected during the run,
does reducing L1D cache size
reduce area, and how strong
is the evidence?''
\end{quote}

To answer this question using C2,
the Tutor must find the relevant
claims and evidence across
the completed run.

It must then determine which
iterations support the relationship,
which provide insufficient evidence,
and whether any comparisons
isolate L1D cache size.

Although the required observations
are represented in C2,
the overall relationship assessment
is not explicitly provided.

This is the limitation
addressed by C3.

% =====================================================
% C3 — ACCUMULATED RELATIONSHIP
% =====================================================

\subsection{C3: Accumulated Relationship Assessment}

C3 extends the Pure Graph
by adding structured assessments
of architectural relationships
observed across multiple iterations.

These assessments are constructed
deterministically from the recorded
runtime evidence.

Therefore, C3 does not introduce
new simulation measurements.
Instead, it summarizes what
the existing observations indicate
about a particular relationship.

For the same L1D-size--area
prediction, C3 provides
the accumulated relationship
shown in
Listing~\ref{lst:app_p5_c3_relationship}.

\begin{lstlisting}[
    style=remodeljson,
    caption={Accumulated L1D-size--area
    relationship assessment in C3.},
    label={lst:app_p5_c3_relationship}
]
{
  "relationship_state": "supported",
  "causal_strength": "limited_isolation",
  "iteration_outcomes": {
    "2": "supported",
    "3": "insufficient_evidence",
    "4": "supported",
    "5": "supported"
  },
  "num_independent_runtime_iterations": 3,
  "evidence_quality": {
    "exact_parameter_isolations": 1,
    "confounded_observations": 6
  }
}
\end{lstlisting}

\paragraph{Understanding the C3 Assessment.}

The field
\texttt{relationship\_state}
is marked as \texttt{supported}.

This means that the accumulated
runtime evidence supports
the predicted relationship
between decreasing L1D cache size
and decreasing area.

The field
\texttt{iteration\_outcomes}
shows where that evidence
came from.

Iterations 2, 4, and 5
support the relationship.

Iteration 3 is marked
as \texttt{insufficient\_evidence},
meaning that it does not
provide enough usable evidence
to determine the relationship.

Thus, the assessment identifies
three supporting runtime iterations.

The field
\texttt{exact\_parameter\_isolations}
has a value of 1.

This indicates that one
recorded comparison changed
only the parameter of interest,
allowing its observed metric
response to be examined
without simultaneous changes
in other architectural parameters.

However, the assessment also
contains six confounded observations,
where other architectural parameters
changed at the same time.

For this reason,
\texttt{causal\_strength}
is classified as
\texttt{limited\_isolation}.

In simple terms, the evidence
supports the predicted relationship,
but the available comparisons
do not justify treating it
as an unrestricted causal rule.

% =====================================================
% C2 VS C3
% =====================================================

\paragraph{What C3 Improves over C2}

The main difference between
C2 and C3 is how the accumulated
evidence is presented to the Tutor.

Under C2, the Tutor receives
individual evidence records
connected through graph edges.

It must find the relevant
observations and combine them
to determine the overall
relationship.

Under C3, this accumulated
assessment is already available
as a structured relationship node.

For example, C3 explicitly
provides that the L1D-size--area
relationship is supported by
three independent runtime
iterations and has
\texttt{limited\_isolation}
causal strength.

The Tutor can therefore use
this accumulated assessment
without independently
reconstructing it from
individual evidence records.

Table~\ref{tab:app_p5_c2_c3_comparison}
summarizes the difference.

\begin{table}[htbp]
\centering
\caption{Evidence available under
C2 and C3 for the same relationship.}
\label{tab:app_p5_c2_c3_comparison}

\renewcommand{\arraystretch}{1.2}
\small
\setlength{\tabcolsep}{5pt}

\begin{tabular}{p{0.48\linewidth}cc}
\toprule
\textbf{Information} &
\textbf{C2} &
\textbf{C3} \\
\midrule

Individual runtime observations &
Yes & Yes \\

Connected claims and evidence &
Yes & Yes \\

Explicit accumulated relationship state &
No & Yes \\

Explicit iteration-level evidence summary &
No & Yes \\

Explicit accumulated causal-strength assessment &
No & Yes \\

\bottomrule
\end{tabular}
\end{table}

Importantly, C3 does not
replace the individual evidence.
It preserves that evidence
and adds an accumulated
assessment derived from it.

Therefore, C3 improves
the information directly
available to the Tutor
for questions requiring
evidence from multiple iterations.

% =====================================================
% C3 LIMITATION
% =====================================================
\paragraph{Remaining Limitation of C3}

Although C3 makes accumulated
relationship assessments explicit,
it still supplies the complete
Evidence-Enriched Graph
to the Tutor.

The graph contains information
about many executed designs,
architectural parameters,
simulation metrics, claims,
evidence records, and lessons.

However, a student usually asks
about only a small part
of this information.

For example, the question
about L1D cache size and area
requires the relevant
L1D-size--area relationship
and its supporting evidence.

It does not require every
other architectural relationship
or observation stored
in the completed run.

Under C3, the Tutor must
identify the relevant
relationship within
the larger graph context.

Therefore, even though
the accumulated assessment
is available, it may be
surrounded by substantial
information unrelated
to the current question.

This is an important distinction:
C3 makes relationship evidence
explicit, but does not perform
question-focused retrieval
before supplying the context
to the Tutor.

% =====================================================
% SUMMARY
% =====================================================

\paragraph{Summary and Motivation for C4.}

This example demonstrates
the additional capability
introduced by C3.

C2 organizes individual
runtime observations
through graph connections.

C3 extends this representation
with accumulated relationship
assessments, allowing the Tutor
to access the overall evidence
state and causal-strength
classification directly.

However, C3 still supplies
the complete enriched graph,
including information
unrelated to the question.

This remaining limitation
motivates C4 Question-Focused
GraphRAG, which retrieves
the relevant evidence
and relationship assessments
from the enriched graph
according to the student's
architectural question.


\section{From Evidence-Enriched Graph to Question-Focused GraphRAG}
\label{app:p5_c4_example}

This section continues the same
Matrix Multiplication example
introduced in
Section~\ref{app:p5_c1_c2_example}
and examined further in
Section~\ref{app:p5_c2_c3_comparison}.

The architectural relationship is:

\begin{quote}
``Reducing \texttt{l1d\_size}
may lower area.''
\end{quote}

The previous section demonstrated
that C3 provides an explicit
assessment of this relationship
using evidence accumulated
across multiple DSE iterations.

However, C3 supplies the complete
Evidence-Enriched Graph to the Tutor,
including information that may
not be relevant to the question.

C4 Question-Focused GraphRAG
addresses this limitation by
retrieving the relevant portion
of the enriched graph.

The purpose of this example
is to demonstrate how C4
selects the required evidence
while preserving the accumulated
relationship assessment.

% =====================================================
% C3 LIMITATION
% =====================================================

\subsection{C3: Complete Evidence-Enriched Graph}

Under C3, the Tutor receives
the complete enriched graph
constructed from the recorded
DSE execution.

The graph includes executed
designs, simulation measurements,
architectural claims, evidence
records, lessons, and accumulated
relationship assessments.

For the L1D-size--area relationship,
C3 explicitly provides:

\begin{itemize}
    \item Relationship state:
    \texttt{supported}.

    \item Supporting iterations:
    2, 4, and 5.

    \item Causal strength:
    \texttt{limited\_isolation}.

    \item Exact parameter-isolated
    comparisons: 1.

    \item Confounded observations: 6.
\end{itemize}

These values allow the Tutor
to identify the accumulated
evidence supporting the
architectural relationship.

However, the complete graph
also contains information
about other architectural
parameters and metrics.

Therefore, the Tutor must
identify which relationship
and evidence records are
relevant to the student's
specific question.

This is the limitation
addressed by C4.

% =====================================================
% QUESTION
% =====================================================

\subsection{Recorded Architectural Question}

The following question was
selected from the recorded
Matrix Multiplication run
\texttt{run\_20260731\_195619}:

\begin{quote}
``Based on all the evidence accumulated
in this run, what relationship was
observed between \texttt{l1d\_size}
and area?

Specifically, what does the runtime
evidence show about whether a decrease
in \texttt{l1d\_size} is associated
with a decrease in area?

How strong is the evidence,
and can this be treated as
an established causal relationship?''
\end{quote}

To answer this question,
the Tutor needs information
about the L1D-size--area
relationship, the supporting
runtime observations, and
the strength of the evidence.

It does not need every other
architectural relationship
stored in the completed run.

% =====================================================
% C4 RETRIEVAL
% =====================================================

\subsection{C4: Question-Focused Evidence Retrieval}

C4 examines the architectural
question to identify the
relevant parameter, metric,
and relationship.

For this example,
the retrieval mechanism
identifies:

\begin{itemize}
    \item Architectural parameter:
    \texttt{l1d\_size}.

    \item Target metric:
    \texttt{area}.

    \item Parameter direction:
    decrease.

    \item Expected metric direction:
    decrease.
\end{itemize}

The retrieval plan uses
these question-specific
anchors to select the
corresponding relationship
and related runtime evidence
from the Evidence-Enriched Graph.

Listing~\ref{lst:app_p5_c4_retrieval}
presents selected fields
from the recorded retrieval plan.

\begin{lstlisting}[
    style=remodeljson,
    caption={Selected C4 retrieval settings
    for the L1D-size--area question.},
    label={lst:app_p5_c4_retrieval}
]
{
  "question_type": "accumulated_relationship",
  "iteration_anchors": [],
  "parameter_anchors": [
    "l1d_size"
  ],
  "metric_anchors": [
    "area",
    "l1d_size"
  ]
}
\end{lstlisting}

Unlike the complete C3 context,
the C4 context is constructed
specifically for the
architectural question.

The retrieval mechanism
selects the relevant
relationship and its
connected evidence rather
than supplying every
relationship in the graph.

% =====================================================
% RETRIEVED RELATIONSHIP
% =====================================================

\subsection{Retrieved Accumulated Relationship}

The retrieved C4 context
contains the same
L1D-size--area relationship
assessment available in C3.

Listing~\ref{lst:app_p5_c4_relationship}
presents selected fields
from the retrieved relationship.

\begin{lstlisting}[
    style=remodeljson,
    caption={Accumulated relationship
    retrieved by C4 GraphRAG.},
    label={lst:app_p5_c4_relationship}
]
{
  "relationship_state": "supported",
  "causal_strength": "limited_isolation",
  "iteration_outcomes": {
    "2": "supported",
    "3": "insufficient_evidence",
    "4": "supported",
    "5": "supported"
  },
  "num_independent_runtime_iterations": 3,
  "evidence_quality": {
    "exact_parameter_isolations": 1,
    "confounded_observations": 6
  }
}
\end{lstlisting}

These values match the
accumulated relationship
assessment presented in C3.

The relationship is classified
as \texttt{supported}, with
supporting evidence from
Iterations 2, 4, and 5.

Iteration 3 provides
insufficient evidence.

The assessment also records
one exact parameter-isolated
comparison and six
confounded observations.

Therefore, the causal
strength is classified as
\texttt{limited\_isolation}.

Importantly, C4 does not
generate a new relationship
assessment or modify the
recorded simulation evidence.

Instead, it retrieves
the existing assessment
from the enriched graph.

% =====================================================
% C3 VS C4
% =====================================================

\subsection{What C4 Improves over C3}

The main difference between
C3 and C4 is how the Tutor
receives the relevant evidence.

Under C3, the Tutor receives
the complete enriched graph
and must identify the
L1D-size--area relationship
within the larger context.

Under C4, the retrieval
mechanism selects that
relationship before the
Tutor generates its answer.

For the same question,
both conditions have access
to the same accumulated
relationship assessment.

However, C4 provides
the relevant assessment
within a smaller,
question-focused context.

Table~\ref{tab:app_p5_c3_c4_comparison}
summarizes the difference.

\begin{table}[htbp]
\centering
\caption{Comparison of evidence access
under C3 and C4.}
\label{tab:app_p5_c3_c4_comparison}

\renewcommand{\arraystretch}{1.2}
\small
\setlength{\tabcolsep}{5pt}

\begin{tabular}{p{0.48\linewidth}cc}
\toprule
\textbf{Characteristic} &
\textbf{C3} &
\textbf{C4} \\
\midrule

Individual runtime evidence &
Yes & Yes \\

Accumulated relationship assessment &
Yes & Yes \\

Complete enriched graph supplied &
Yes & No \\

Question-focused retrieval &
No & Yes \\

\bottomrule
\end{tabular}
\end{table}

This comparison shows
that C4 does not improve
the underlying simulation
evidence itself.

Instead, C4 improves
how the relevant evidence
is selected and supplied
to the Tutor.

% =====================================================
% INTERPRETATION
% =====================================================

\subsection{Interpretation of the C3--C4 Comparison}

The example demonstrates
the role of question-focused
retrieval in architectural
question answering.

C3 provides the accumulated
relationship assessment,
but the Tutor must locate
it within the complete graph.

C4 retrieves the relevant
relationship and its
associated evidence,
reducing the amount of
unrelated information
supplied to the Tutor.

This distinction is
consistent with the
aggregate Part B results,
where C4 achieves higher
evidence utilization
while requiring
substantially less context.

However, the retrieved
evidence must still
be interpreted correctly.

Question-focused retrieval
improves evidence selection,
but does not guarantee
that every generated
architectural explanation
will be factually correct.

\paragraph{Summary of the C1--C4 Progression.}

Together, the three
case studies illustrate
the progression of the
four evidence conditions.

C1 preserves the
simulation-grounded
runtime experiences.

C2 organizes those
experiences through
explicit graph connections.

C3 adds accumulated
relationship assessments
derived from the
recorded evidence.

C4 retrieves the
question-relevant portion
of the enriched graph
before generating
the architectural explanation.

The progression demonstrates
how evidence preservation,
organization, enrichment,
and retrieval address
different challenges
in simulation-grounded
architectural question answering.


\section{Comparison of Recorded Tutor Responses}
\label{app:p5_comparative_case_study}

To examine how the four evidence conditions
affect architectural explanations, we compare
their recorded responses to the same question
from Matrix Multiplication run
\texttt{run\_20260731\_195619}.

The question asks whether decreasing
\texttt{l1d\_size} is associated with
decreasing area, how strong the accumulated
evidence is, and whether the relationship
can be considered causal.

Table~\ref{tab:app_p5_response_comparison}
presents selected verbatim excerpts
from the recorded Tutor responses,
together with their strengths
and observed limitations.

\begin{table*}[t]
\centering
\caption{Comparison of recorded Tutor
explanations for the same
L1D-size--area relationship question.}
\label{tab:app_p5_response_comparison}

\renewcommand{\arraystretch}{1.3}
\footnotesize
\setlength{\tabcolsep}{4pt}

\begin{tabular}{
p{0.08\textwidth}
p{0.38\textwidth}
p{0.23\textwidth}
p{0.23\textwidth}
}
\toprule
\textbf{Condition} &
\textbf{Recorded Response Excerpt} &
\textbf{Strength} &
\textbf{Limitation / Observation} \\
\midrule

C1 &
``The evidence strength is moderate;
while there is supporting evidence
for the claim that reducing
l1d\_size lowers area, the presence
of confounding parameters in the
comparisons means that a strong
causal relationship cannot be
firmly established.'' &
Identifies the association and
recognizes that other parameters
may affect area. &
Does not report the complete
accumulated relationship assessment.
Treats area variation at unchanged
L1D size as evidence of an
inconsistent relationship. \\

\addlinespace[5pt]

C2 &
``The evidence strength is moderate,
as there are several supporting
observations, but the exact
parameter isolation is false due
to the simultaneous changes in
other parameters.'' &
Identifies multiple supporting
observations and recognizes
confounding. &
Does not report
\texttt{limited\_isolation},
despite one exact isolated
comparison in the accumulated
evidence. \\

\addlinespace[5pt]

C3 &
``This relationship has been
supported by multiple iterations,
with three instances confirming
the association, while two
instances were classified as
insufficient evidence due to
confounding factors.'' &
Uses accumulated relationship
information and reports
supporting evidence counts. &
Does not clearly distinguish
claim-evidence counts from
independent iteration outcomes;
the stated reason for the
insufficient-evidence records
is not fully established. \\

\addlinespace[5pt]

C4 &
``This relationship is supported
by evidence from iterations 2,
4, and 5, where decreases in
l1d\_size corresponded with
decreases in area.'' &
Identifies the supporting
iterations and provides
the relevant numerical
comparisons and causal strength. &
Does not explicitly distinguish
the two insufficient-evidence
records from independent
iteration outcomes. \\

\bottomrule
\end{tabular}
\end{table*}

The comparison shows that all four
conditions identify the general
association between decreasing
L1D cache size and decreasing area.

However, the completeness and
precision of the explanations differ.

C1 primarily uses individual
chronological observations,
while C2 identifies supporting
evidence across multiple iterations.

C3 additionally reports
accumulated relationship
information, including evidence
counts and causal strength.

C4 provides a question-focused
explanation that identifies the
relevant supporting iterations
and their numerical comparisons.

These observations illustrate
how evidence organization,
enrichment, and retrieval
influence the information
incorporated into generated
architectural explanations.

The comparison concerns one
recorded question and complements
the aggregate Evidence Utilization,
Claim Correctness, and
Question-Answering Quality
results presented in Part B.


% =====================================================
% APPENDIX — LLM VERIFIER DISAGREEMENT
% =====================================================

\section{LLM Verifier Disagreement Case Study}
\label{app:part_b_verifier_disagreements}

This case study illustrates a
disagreement between automated
LLM-based verification and
human evaluation of an
architectural explanation.

The example is drawn from
the O3CPU scenario,
run \texttt{run\_20260820\_001258},
Question Q3, under C3
Evidence-Enriched Graph.

The question examines architectural
parameter changes between
Iterations 4 and 5 and whether
the observed IPC improvement
can be attributed to a
specific parameter.

\subsection{Original Architectural Claim}

The recorded Tutor response
included the following claim:

\begin{quote}
``\texttt{num\_iq\_entries}
increased from 16 to 32.''
\end{quote}

This statement predicts a
specific architectural transition
between two executed configurations.

Because it identifies a parameter
and exact numerical values,
its correctness can be checked
directly against the recorded
design configurations.

\subsection{Recorded Simulation Evidence}

The authoritative execution
records show that
\texttt{num\_iq\_entries}
remained unchanged between
Iterations 4 and 5.

Table~\ref{tab:app_verifier_disagreement}
summarizes the claim,
recorded evidence,
and evaluation judgments.

\begin{table}[htbp]
\centering
\caption{Recorded disagreement between
LLM verification and human evaluation
for an O3CPU architectural claim.}
\label{tab:app_verifier_disagreement}

\renewcommand{\arraystretch}{1.2}
\small
\setlength{\tabcolsep}{5pt}

\begin{tabular}{p{0.55\linewidth}p{0.30\linewidth}}
\toprule
\textbf{Measurement} &
\textbf{Recorded Value} \\
\midrule

Claim identifier & CL3 \\

Claimed IQ entries &
16 $\rightarrow$ 32 \\

Actual IQ entries &
32 $\rightarrow$ 32 \\

Actual parameter change &
0 \\

Deterministic evidence status &
\texttt{CONTRADICTED} \\

LLM verifier correctness score &
1.0 \\

LLM verifier confidence &
95\% \\

Human correctness score &
0.0 \\

\bottomrule
\end{tabular}
\end{table}

The claim states that the
number of instruction queue
entries increased.

However, both executed
configurations contain
32 instruction queue entries.

Therefore, no increase
occurred during the
specified transition.

The deterministic checker
correctly identifies this
discrepancy and classifies
the claim as
\texttt{CONTRADICTED}.

\subsection{Automated Verifier Judgment}

Despite the contradictory
simulation evidence,
the LLM verifier assigned
the claim a correctness
score of 1.0 with
95\% confidence.

Its recorded explanation was:

\begin{quote}
``The claim that num\_iq\_entries
increased from 16 to 32 is directly
supported by the design records
for iteration 5.''
\end{quote}

This judgment is incorrect
because verifying an
architectural transition
requires comparing both
the previous and current
parameter values.

The Iteration 5 value
alone does not establish
that an increase occurred.

The actual comparison
shows that the parameter
remained unchanged.

\subsection{Human Evaluation and Interpretation}

The human evaluator assigned
the same claim a correctness
score of 0.0.

This judgment agrees with
the deterministic evidence
checker and the recorded
architectural configurations.

The disagreement illustrates
a specific limitation of
automated LLM verification:
the verifier accepted an
incorrect parameter-transition
claim despite the availability
of contradictory numerical
evidence.

Importantly, the verifier
reported 95\% confidence
in its incorrect judgment.

The same response also
received a complete-response
Question-Answering Quality
score of 1.0 from the
LLM verifier, compared
with 0.5 under human
evaluation.

Thus, the error affected
both the claim-level
assessment and the
overall evaluation of
the architectural explanation.

\paragraph{Implications for Automated Verification.}

This example demonstrates
that high verifier confidence
does not guarantee agreement
with authoritative simulation
evidence or human evaluation.

Even a simple numerical
parameter transition can
be incorrectly accepted
when the verifier does
not adequately check
both endpoints.

The case supports the
aggregate Part B findings
that automated verification
may overestimate architectural
explanation correctness.

It also highlights the
importance of combining
LLM-based evaluation with
deterministic numerical
checks and independent
human review.

Although this case illustrates
one recorded failure, the
broader reliability of the
automated verifier is
evaluated through the
aggregate claim-level,
response-level, and
confidence-calibration
analyses presented in Part B.
